% Chapter 5: Conclusion and Future Directions

\chapter{Conclusion and Future Directions}
\label{chap:conclusion}

\section{Summary of Research Findings}
This thesis addressed the compounding challenges of the \textbf{Tokenizer Tax} and \textbf{Over-Provisioning} in Large Language Model inference for Hindi-English (Hinglish) code-mixed discourse. Through the design, implementation, and empirical validation of an OpenAI-compatible 12-module serving gateway, we established the following core findings:

\begin{enumerate}
    \item \textbf{Transferability of Prompt Compression (RQ1):} At achieved kept token ratios between $0.50$ and $0.70$, extractive prompt compression preserves conversational task success (intent recognition and slot-filling accuracy) within a pre-registered $-0.05$ non-inferiority margin, demonstrating that nearly half of input prompt tokens can be safely pruned on CPU.
    \item \textbf{Economic Viability (RQ2):} The combination of prompt compression and cascade routing yields a 63.1\% net financial reduction on production inference workloads, mitigating the 2.14$\times$ tokenizer penalty suffered by multilingual users under standard Western subword vocabularies.
    \item \textbf{Identification of Failure Modes (RQ3):} We proved that accuracy failures in unguided neural prompt compression are heavily concentrated in structured entities, timestamps, phone numbers, and monetary figures, where raw classifiers exhibit a 20--31\% error rate.
    \item \textbf{Provable Entity Preservation (RQ4):} By introducing regular-expression safety gates and fail-closed reinjection (\texttt{[[PSi]]} masking), we proved that entity recall can be restored to exactly \textbf{1.000} (100\%) byte-for-byte without compromising compression rates.
    \item \textbf{Statistical Cascade Routing:} Deploying a fine-tuned ModernBERT difficulty scorer under Hoeffding Learn-Then-Test conformal prediction allows operators to bound cheap model routing failures below explicit risk budgets ($\alpha = 0.05$).
\end{enumerate}

\section{Practical Deployment Insights}
From an engineering perspective, several actionable principles emerged from building this gateway:
\begin{itemize}
    \item \textbf{CPU-Only Pre-Processing is Essential:} For a gateway to remain economically viable, prompt compression must execute in $<10$\,ms on commodity CPU hardware without consuming GPU clusters. Our Linear-Chain CRF and heuristic filters fulfill this criterion.
    \item \textbf{Fail-Closed Design Prevents Catastrophic Errors:} In production customer service and transactional systems, an uncompressed prompt is infinitely preferable to a corrupted prompt with broken phone numbers or amounts.
    \item \textbf{Auditability Drives User Trust:} Generating cryptographic usage receipts for every query provides users and enterprise operators with auditable visibility into token savings, model selection, and risk bounds.
\end{itemize}

\section{Limitations}
We explicitly acknowledge the following constraints in our current experimental evaluation:
\begin{itemize}
    \item \textbf{Corpus Script Representation:} While the PRESTO benchmark contains mixed Devanagari and Latin-script task interactions, Latin-script Romanized Hinglish is characterized by extreme orthographic variance that demands continuous expansion of the phonetic lexicon.
    \item \textbf{Evaluator API Rate Limits:} Large-scale automated evaluation over thousands of dev-split dialogues encountered upstream API rate limits (HTTP 429), necessitating batching, retries, and offline caching.
    \item \textbf{Model Proxies:} In the free-tier experimental tracks, high-capacity models such as Llama-3.3-70B served as empirical proxies for closed frontier models.
\end{itemize}

\section{Directions for Future Work}
Several promising avenues for continued exploration remain:
\begin{enumerate}
    \item \textbf{Edge and Mobile Deployment:} Porting the lightweight Linear-Chain CRF and phonetic normalization engine to Android and iOS mobile keyboard runtimes (e.g., via ONNX or WebAssembly), compressing prompts directly on the user's device prior to network transmission.
    \item \textbf{Online Conformal Risk Adaptation:} Transitioning from static offline calibration splits to online, streaming conformal prediction algorithms that adjust the routing threshold $\tau^*$ dynamically in response to non-stationary concept drift in incoming user queries.
    \item \textbf{Expansion to Dravidian and Indo-Aryan Languages:} Extending the phonetic consonant-skeleton algorithm and code-mixing lexicons to other major Indian language combinations, including Tanglish (Tamil-English), Tenglish (Telugu-English), and Banglish (Bengali-English).
\end{enumerate}
